\documentclass[10pt,a4paper]{amsart}
\usepackage[margin=19mm]{geometry}
\usepackage{amsmath,amssymb,mathtools}
\usepackage[scr=rsfs,cal=euler]{mathalfa}
\usepackage{xfrac}
\usepackage[unicode,hidelinks,bookmarks=false]{hyperref}
\newcommand{\ie}{i.e.\ }
\newcommand{\cf}{cf.\ }
\newcommand{\resp}{respectively\ }
\newcommand{\Subsection}{\subsection}
\providecommand{\nobreakdash}{\nobreak\hbox{-}}
\setlength{\emergencystretch}{2em}
\multlinegap0pt
\usepackage{csquotes}
\usepackage{amssymb}
\usepackage{mathtools}
\usepackage{xcolor}
\usepackage{tikz}
\usepackage{float}
\newtheorem{lemma}{Lemma}
\newcommand{\retract}{\mathbin{\rotatebox[origin=c]{-45}{$\rightsquigarrow$}}}
\newcommand{\s}{\mathbb{S}}
\title{The fundamental quandle of ribbon concordances}
\author{Eva Horvat, Luka Mar\v{c}i\v{c}}
\date{Source: arXiv:2602.20841v2 (CC BY 4.0)}
\begin{document}
\maketitle

\noindent\emph{Source and licence.} The fundamental quandle of ribbon concordances. Version arXiv:2602.20841v2 (CC BY 4.0), distributed by arXiv under the Creative Commons Attribution 4.0 International licence, http://creativecommons.org/licenses/by/4.0/, as recorded in the arXiv licence metadata for this exact version and shown on the abstract page https://arxiv.org/abs/2602.20841v2. Full licence notice: \texttt{licenses/arxiv-2602.20841-CC-BY-4.0.txt}.\par\smallskip

\noindent\textit{Editorial scope.} Counted unit: the article's lemma lemma:ribon\_conc\_quandle (source0.tex lines 734--742). The article's complete original proof of it is delivered with it (source0.tex lines 743--785, one \texttt{proof} environment of 43 lines). No mathematics is written, repaired or re-derived: the statement and the proof are the article's own lines, in the article's own notation, and no retained line is substituted or re-flowed.  Retained uncounted as the source-local context the proof needs: lines 788--794; lines 797--804.  One declared adaptation: exactly 2 ancillary strings inside the retained spans are omitted (image-inclusion commands and, where the source repeats a label, the repeated label definitions) and no other line differs from the source; the figure environments, with their placement, captions and labels, are kept, and the package records the omitted strings in fidelity receipts bound to the affected bodies.  No retained line contains a citation, so the wrapper carries no bibliography and no bibliography entry.  The retained lines contain the article's own diagram code, which the wrapper typesets with the standard tikz packages; no image file is delivered and the package declares no asset dependency.  Editorial formatting only, in addition: an AMS \texttt{amsart} preamble that restates the article's own declarations and macros used by the retained lines, namely the environments \texttt{lemma} on separate counters, together with the standard packages those lines need.  The retained environments print in this document as Lemma 1 in order of appearance, whereas the article prints the same items with the numbering of its own full document.  The complete native source of the article as distributed in the arXiv e-print for this exact version is bundled unchanged as \texttt{sources/195224/source0.tex}.
\begin{figure}[ht]
	\centering
	%\includesvg[width =0.6\textwidth]{svg-inkscape/ribbon_concordance.svg}
	
	\caption{Example of a ribbon concordance with one critical point of index 0 and 1 each.}
	\label{fig:ribbon_conc}
\end{figure}

\begin{figure}[ht]
	\centering
	\hspace{1cm}
	%\includesvg[width =0.45\textwidth]{svg-inkscape/concordance_neighbourhood.svg}
	
	\caption{Fixing $H(0,s)$ so that it doesn't pass $N_C((0,1)_D \times \{1\}_D)$.}
	\label{fig:conc_nbh}
\end{figure}

\begin{lemma}
	\label{lemma:ribon_conc_quandle}
	Let $C_0\subset \s ^{3}\times [-1, 0]$ be a ribbon concordance from a knot $K_{0}$ to a knot $K_{-1}$,
	and let $C_1 \subset \s^3 \times [0,1]$ be a ribbon concordance from a knot $K_1$ to knot $K_0$, 
	where $K_{i}\in \s ^{3}\times \{i\}$.
	Assume the projection of $C_0$ and $C_1$ to $[-1,1]$ is a Morse function and that the projection of $C_1$ to $[0,1]$ has a single critical point of index 0 and a single
	critical point of index 1. Let $C = C_0 \cup C_1$ be a ribbon concordance from knot $K_1$ to knot $K_{-1}$.
	Then the quandle homomorphism $\iota \colon Q(C_0)\to Q(C)$, induced by inclusion, is injective.
\end{lemma}

\begin{proof} 
	We begin by establishing some notation. 
	We may assume that the ribbon concordance $C_1$ ends at the exact moment when the 1-handle has been attached (it doesn't meaningfully change after that).
	Then $C_1$ consists of a part that is ambient isotopic to $K_0 \times I$ (so that is what we will call it and regard it as), a disk $D$
	and a band $B$, connecting $K_0 \times I$ to $D$ (see Figure~\ref{fig:ribbon_conc} for an example). Define normal disk bundles
	$N_{C_0} \subset \s^3 \times [-1,0]$ and $N_C \subset \s^3 \times [-1,1]$ about $C_0$ and $C$, respectively,
	along with their respective exteriors $E_{C_0} \subset \s^3 \times [-1,0]$ and $E_C \subset \s^3 \times [-1,1]$.
	Notice $N_{C_0}  \subset N_C$ and $E_{C_0} \subset E_C$.
	Lastly, choose a basepoint $z \in E_{C_0}$ for which we define fundamental quandles $Q(C_0)$ and $Q(C)$.


	We can regard both disk $D$ and band $B$ as the product $I \times I$. For clarity's sake we write $D = I_D \times I_D$ and $B = I_B \times I_B$, where $B$ is attached
	to some arc of $K_0 \times {1}$ {along} $I_B \times \{0\}$, and is attached to $I_D \times \{0\} \subset D$ {along} $I_B \times \{1\}$ in the natural way.
	We have a series of strong deformation retractions
	{ \small
	$$C_1 = (K_0 \times I) \cup B \cup D \overset{I_D \times I_D \retract I_D \times \{0\}}{\longrightarrow} (K_0 \times I) \cup B 
	\overset{I_B \times I_B \retract I_B \times \{0\}}{\longrightarrow} K_0 \times I
	\overset{K_0 \times I \retract K_0 \times \{0\}}{\longrightarrow} K_0$$
	}
	which, in turn, give us a strong deformation retraction $R \colon N_C \times I \to N_{C_0}$.


	Since multiple copies of the interval $I$ with distinct meanings will appear from now on, we index them by the variable we use for them; $t$ for $I_t$, $s$ for $I_s$ and $r$ for $I_r$.
	Let $[\alpha], [\beta] \in Q(C_0)$ be such that $\iota([\alpha]) = \iota([\beta]) \in Q(C)$. By definition, $\alpha$ and $\beta$ are paths 
	$(I_t, \{0\}_t, \{1\}_t) \to (E_{C_0}, \partial N_{C_0}, \{z\})$, and we have a homotopy $H\colon I_t \times I_s \to E_{C}$ such that $H(t,0) = \alpha(t)$, $H(t,1) = \beta(t)$,
	$H(0,s) \in \partial N_C$ and $H(1,s) = z$.
	Fixing $t=0$, $H(0,s)$ defines a path in $\partial N_{C}$ from $\alpha(0)$ to $\beta(0)$. The retraction $R \colon N_C \times I_r \to N_{C_0}$ then defines 
	a homotopy $\widetilde{H}_0\colon \{0\}_t \times I_s \times I_r \to \partial N_{C} \subset E_C$,
	where $\widetilde{H}_0(0, s, 0) =  H(0,s)$ and $\widetilde{H}_0(0,s,1)$ is some path in $\partial N_{C_0}$ from $\alpha(0)$ to $\beta(0)$.
	Note that, a priori, the retraction $R$ need not map $H(0,s)$ entirely to $\partial N_C$ for every $r$; it may happen that $R(H(0,s),r) \in \text{int}(N_C)$ for some $s$ and $r$!
	It is clear from the relatively simple definition of $R$ that this happens exactly when $H(0,s)$ moves over a part of $\partial N_C$ defined by the boundary of the disk bundle of some point in $(0,1)_D \times \{1\}_D$.
	To avoid this, we can compose $R$ with a homotopy that first moves $H(0,s)$ along $\partial N_C$ 
	into a favourable position (outside the disk bundles of points in $(0,1)_D \times \{1\}_D$) where this doesn't happen, thus obtaining $\widetilde{H}_0$ as desired
	(see figure~\ref{fig:conc_nbh}).
	Define also $\widetilde{H}_0(1,s,r) = z$.
	Now, since $(I_t \times I_s, \{0,1\}_t \times I_s )$ is a CW-pair, it has the homotopy extension property. 
	Thus, the homotopy $\widetilde{H}_0$ can be
	extended to a homotopy $\widetilde{H}\colon I_t \times I_s \times I_r \to E_C$ such that 
	$\widetilde{H}(0,s,r) = \widetilde{H}_0(0,s,r) \in  \partial N_C, \widetilde{H}(1,s,r) = \widehat{H}_0(1,s,r) = z$ and $\widetilde{H}(t,s,0) = H(t,s)$.
	As such, $\widetilde{H}(t,s,1)$ is then a homotopy in $E_C$ between $\alpha$ and $\beta$ such that $\widetilde{H}(0,s,1) \subset \partial N_{C_0}$ and $\widetilde{H}(1,s,1) = z$.
	Projecting $\widetilde{H}(t,s,1)$ onto $E_{C_0}$ then gives us the desired homotopy between $\alpha$ and $\beta$,
	showing that $[\alpha] = [\beta]$ as elements of $Q(C_0)$.
\end{proof}

\end{document}
